\chapter{Global Anomalies, Spontaneous Symmetry Breaking, and Pions}
\label{ch:volume-ii-21}

\section{Anomalies Of Global Symmetries}

A global symmetry may have an anomaly when coupled to background gauge fields.
This does not make the theory inconsistent. Instead it says that the symmetry
cannot be gauged without additional degrees of freedom or inflow. Such an
anomaly is invariant under RG flow. The infrared theory must reproduce it.

\section{Spontaneous Symmetry Breaking}

In an infinite-volume quantum system, a continuous global symmetry \(G\) may be
spontaneously broken to \(H\). The vacuum is invariant under \(H\), and the
low-energy fields take values in the coset
\[
  G/H.
\]
Goldstone modes are coordinates on this manifold. Their interactions are
derivative at leading order because constant shifts along the vacuum manifold
cost no energy.

\section{Chiral Symmetry In QCD}

For \(N_f\) massless quark flavors, the non-anomalous chiral symmetry is
approximately
\[
  SU(N_f)_L\times SU(N_f)_R
  \longrightarrow
  SU(N_f)_{\rm vector}.
\]
The Goldstone fields are assembled into
\[
  U(x)\in SU(N_f),
\]
with leading effective Lagrangian
\[
  \mathcal L_{\pi}
  =
  \frac{f_\pi^2}{4}\operatorname{Tr}
  (\partial_\mu U^\dagger\partial^\mu U)+\cdots.
\]

\section{Wess--Zumino--Witten Terms}

Anomaly matching requires terms in the pion effective action not visible from
ordinary derivative counting alone. The Wess--Zumino--Witten term reproduces
the chiral anomaly of the microscopic quarks. It explains, for example, the
anomalous coupling responsible for \(\pi^0\to\gamma\gamma\). The low-energy
theory remembers the ultraviolet anomaly exactly.
